\documentclass[12pt]{article}
\usepackage[utf8]{inputenc}
\usepackage[T1]{fontenc}
\usepackage{hyperref}
\title{Francisco Le\'on Zoque Language Resources for K-12 Curriculum Development}
\author{Compiled by Claude Sonnet 5.5 (claude-sonnet-5-5) for Andr\'as Kornai}
\date{}
\begin{document}
\maketitle

\section{Introduction}
Francisco Le\'on Zoque (ISO 639-3 zos) is a Mixe-Zoquean language of Chiapas, Mexico, with roughly 20,000 speakers. It is part of the Chiapas Zoque dialect cluster. It is reportedly about 83\% mutually intelligible with Copainal\'a Zoque. The town of Francisco Le\'on was destroyed by the 1982 El Chich\'on eruption and its residents were resettled. My searches found little material specific to this variety, so this report is short.

\section{Primary Educational Resources}
The Colegio de M\'exico repository holds literacy primers produced by SEP and the Summer Institute of Linguistics for the Francisco Le\'on dialect of Zoque. One example is \emph{Primera cartilla zoque: dialecto de Francisco Le\'on}. The search results did not show which of the three repository records is which, so check each record before relying on it.
\begin{itemize}
\item Colegio de M\'exico repository record (Zoque literacy material): \url{https://repositorio.colmex.mx/concern/books/zk51vj76r?locale=en}
\end{itemize}

\section{Linguistic and Reference Materials}
\begin{itemize}
\item Wikipedia, ``Chiapas Zoque'' (background on the dialect cluster): \url{https://en.wikipedia.org/wiki/Chiapas_Zoque}
\item Wikipedia, ``Zoque languages'': \url{https://en.wikipedia.org/wiki/Zoque_languages}
\end{itemize}
I found no modern dictionary or grammar specific to Francisco Le\'on Zoque. The grammar and vocabulary items that came up in search were for other Zoque varieties, so I do not list them here.

\section{Religious and Community Resources}
\begin{itemize}
\item Global Recordings Network, Zoque, Francisco Le\'on (evangelistic and Bible-teaching audio): \url{https://globalrecordings.net/en/language/zos}
\item Global Recordings Network, Zoque, Francisco Le\'on: Chapultenango (a related variety, with audio and slideshow downloads): \url{https://globalrecordings.net/en/language/18726}
\end{itemize}

\section{Digital and Audio Resources}
The only audio found is the Global Recordings material listed above.

\section{Curriculum Development Context}
The language is not known to be taught in schools. The documented base for a native-language curriculum is thin: old SIL primers, a small amount of religious audio, and encyclopedic descriptions. A curriculum project would probably need new materials developed with the community. Possible starting points are the SIL primers (after checking their records) and the related Chiapas Zoque varieties.

\end{document}